\chapter{Grundlagen und Analyse der Anforderungen}\label{chap:Hintergrund}
In diesem Kapitel werden die Anforderungen an das Verhalten des Verfahrens, die sich aus \chapref{Zielsetzung} ergeben, näher erörtert.\\
Ebenso werden Grundlagen sowie Anforderungen bezüglich der Testausgabeformate vorgestellt. 

\section{Anforderungen an das Verfahren}\label{Anforderungen Tool}
Folgende Anforderungen werden an einen Durchlauf des Verfahrens gestellt:
\begin{itemize}
    \item \textbf{Die Wahl der verwendeten Testframework(s) ist unerheblich}\\
    Das Verfahren soll framework-agnostisch funktionieren, auch wenn innerhalb des Projektes verschiedene Testframeworks zum Einsatz kommen.
    \item \textbf{Nach Abschluss des Verfahrens wurde jeder Testcase für den aktuellen Stand des Codes mindestens einmal durchgeführt}\\
    Testcases können vor Duchlauf des Verfehrens mehrfach ausgeführt worden sein, allerdings darf während des Verfahrens selbst kein Testcase ausgeführt werden, der bereits ausgeführt wurde, um die Laufzeit der Testphase so gering wie möglich zu halten. Trotzdem muss die gesamte Testsuite komplett abgedeckt sein, um die Aussagekraft des Testlaufes im Vergleich zu einem komplett neuen Durchlauf nicht zu beeinträchtigen.
    \item \textbf{Die Ausgabe des durch das Verfahren ausgeführten Testlaufs ist auf dieselbe Weise verfügbar, wie die anderen Ausgaben}\\
    Die Ausgabe soll in demselben Format, auf dieselbe Art und Weise abgelegt werden, wie auch die anderen Ausgaben, die dem Verfahren zur Verfügung gestellt werden.
    \item \textbf{Die Laufzeit eines Testdurchlaufs unter Verwendung des Verfahrens ist kleiner oder gleich der Laufzeit eines konventionellen Durchlaufs}\\
    Das Vefahren soll die Gesamtlaufzeit der Testausführung reduzieren.
\end{itemize}

\section{Anforderungen an ein Testausgabeformat}\label{Anforderungen Testausgabeformat}
Um ein framework-agnostisches Verfahren zu ermöglichen, muss ein testframework-agnostisches Testausgabeformat gefunden werden, in welchem die Ergebnisse eines Testlaufs in einer standartisierten Form festgehalten werden.\\
\texttt{testAuditor} ist der Kern des Verfahrens, und benötigt zur korrekten Funktion Zugriff auf die Gesamtheit aller Testcases, sowie Informationen über den Status von bereits gelaufenen Testcases. Eine detailliertere Beschreibung von \texttt{testAuditor} findet sich in \chapref{Komponente_testAuditor}.\\
% Jedes Testframework verfügt über eines oder meherere mögliche Ausgabeformate, welches die Informationen über das Ergebnis eines Testlaufes enthält.\\
% Damit ein Tool entwickelt werden kann, welches möglichst frameworkunabhägig funktioniert, muss ein Ausgabeformat gefunden werden, dass von möglichst vielen verschiedenen Testframeworks auf dieselbe Art und Weise verwendet wird, damit die Dateien vom Tool über eine Schnittstelle eingelesen werden können.\\
Mit Blick auf diese, sowie die unter \chapref{Anforderungen Tool} beschriebenen Anforderungen an das Verfahren, würde ein ideales Testausgabeformat sämtliche hier gestellten Anforderungen erfüllen:

\textbf{Ein ideales Ausgabeformat für Tests:}
\begin{itemize}
    \item \textbf{ermöglicht die vollständige Ermittlung aller im Projekt vorhandenen Testcases}\\
    Dieser Punkt ist essenziell, zum einen damit \texttt{testAuditor} eine Aussage darüber treffen kann, ob eine vollständige Abdeckung aller Testcases erzielt wurde, zum anderen muss \texttt{testAuditor} zum Erzielen einer vollständigen Abdeckung potentiell auf jeden vorhandenen Testcase zugreifen können. Dazu muss \texttt{testAuditor} selbstverständlich auch jeder Testcase innerhalb des Projektes bekannt sein.
    \item \textbf{identifiziert jeden Testcase eindeutig}\\
    Je nach Testframework und Implementierung können mehrere Testcases denselben Namen besitzen, sofern sie sich in unterschiedlichen Namespaces oder Kontexten befinden. Das Ausgabeformat muss daher eine eindeutige Zuordnung zwischen den einzelnen Testcases und ihren jeweiligen Ergebnissen ermöglichen. 
    \item \textbf{sagt für jeden Testcase einzelnd aus, ob er ausgeführt wurde oder nicht}\\
    Dieser Punkt setzt die Erfüllung der ersten beiden Punkte voraus. Zusätzlich muss aus dem Eintrag eines Testcases im Ausgabeformat erkennbar sein, ob der Test ausgeführt worden ist.
    \item \textbf{ist ein strukturiertes Format}\\
    Es muss eine erfassbare Struktur für das Testausgabeformat vorhanden sein, damit ein entsprechender Parser für \texttt{testAuditor} entwickelt werden kann.
    \item \textbf{ist in vielen Frameworks über viele Programmiersprachen hinweg vertreten}\\
    Es ist kaum realistisch anzunehmen, dass ein strukturiertes Format existiert, welches von jedem existierenden Testframework gleichermaßen benutzt wird. Dennoch gibt es durchaus gängige Formate, die von unterschiedlichen Testframeworks aus unterschiedlichen Sprachen unterstützt werden, wie folgend in \chapref{Vergleich Testausgabefomate} untersucht wird. 
\end{itemize}

\section{Deterministische builds - 'it works on my machine'}\label{chap:determinism}
Damit in dem hier vorgestellten Verfahren eine qualifizierte Aussage über die Testergebnisse getroffen werden kann, muss gewährleistet sein, dass alle Testcases unabhängig vom ausführenden System immer dasselbe Testergebnis liefern. Beispielsweise muss ein Testcase, der auf der Maschine eines Entwicklers ein bestimmtes Ergebnis liefert, auch zwingend auf einem Build-Server dasselbe Ergebnis liefern. Andernfalls könnte eine Diskrepanz zwischen den Testergebnissen zu einer verfälschten Aussage des Verfahrens über den aktuellen Stand der Tests führen. 

Im Rahmen dieser Bachelorarbeit wird hierfür der Linux-package manager \texttt{Nix} benutzt. \texttt{Nix} bietet die Funktionalität, die Entwicklungsumgebung deklarativ zu beschreiben, indem benötigte Abhängigkeiten in einer sogenannten \texttt{flake} festgelegt werden. Eine dazugehörige Datei \texttt{flake.lock} enthält dann detaillierte Informationen über den build. Dadurch können auf verschiedenen Systemen identische Entwicklungsumgebungen aufgebaut werden. 

Es gilt zu beachten, dass durch diesen Aufbau nur explizit deklarierte interne Abhängigkeiten deterministisch erfasst werden können. Externe Faktoren wie etwa Netzwerkschnittstellen, zufallbasierte Ereignisse oder Ähnliches können auf diese Weise nicht reproduzierbar festgehalten werden.\\
Ein vollständige Garantie des Determinismus zwischen Systemen ist nur schwer zu erreichen, jedoch bietet \texttt{Nix} in diesem Fall eine ausreichende Grundlage. 

\section{Entstehung eines Testausgabeformats}
Um zu verdeutlichen, wie die finale Form eines Testausgabeformates entsteht, wird in diesem Kapitel näher auf die Faktoren eingegangen, welche die genaue Struktur und das finale Aussehen des Testausgabeformates beeinflussen.
Die Form eines Testausgabeformates hängt im Wesentlichen von den folgenden Faktoren ab, deren Aufteilung der von Winkler et al.\cite{winkler} vorgestellten Struktur eines Testframeworks folgt: % TODO test discovery hinzufügen?
\begin{itemize}
  \item \textbf{verwendetes Testframework:}\\
  Das verwendete Testframework ist für gewöhnlich an jedem Schritt der Testgenerierung beteiligt und erfült die Rolle eines Managers der einzelnen Komponenten. Daher bestimmt es je nach konkretem Framework stärker oder schwächer die Verwendung der folgenden Punkte:
  \item \textbf{Struktur der zu grunde liegenden Testdateien:}\\
  Die (Nicht-)Verwendung von Klassen, Testsuites oder anderen vom Testframework oder der Programmiersprache bereitgestellten Optionen zur Organisation der Testcases, sowie die Dateistruktur der Testdateien können einen Einfluss auf das Endresultat haben. 
  \item \textbf{Test Runner:}\\
  Auch als \textit{Scheduler, Harness} oder \textit{Orchestrator} bezeichnet, beschreibt der Begriff ''Test Runner'' hier die Komponente eines Testframeworks, welche die Testcases ausführt und die Ergebnisse an den Test Reporter weiterreicht. Der Test Runner beeinflusst u.A die Darstellung von Ereignissen, die während des Testlaufs auftreten und in die Ausgabe übernommen werden.
  \item \textbf{Test Reporter:}\\
  Auch als \textit{test generator} bezeichnet, beschreibt ''Test Reporter'' hier jene Komponente eines Testframeworks, welche aus dem vom Test Runner bestimmten Ergebnis der Testsuite die menschen- bzw. maschinenlesbare Ausgabe generiert.
  \item \textbf{verwendetes Testausgabeformat:}\\
  Damit ist sowohl das Dateiformat, als auch die Wahl eines strukturierten Ausgabeformats gemeint. Näheres dazu in \chapref{Vergleich Testausgabefomate}.
\end{itemize}
In den meisten Testframeworks übernimmt das Testframework selbst die Aufgaben des Test Runners und des Test Reporters, doch in manchen Fällen werden diese Komponenten aber auch durch andere Prozesse bereitgestellt. Dies ist zum Beispiel in der Programmiersprache Java der Fall: Hier stellen die beiden build-tools Maven und Gradle jeweils ihre eigenen Test Runner bzw. Test Reporter zur Verfügung, was unterschiedliche Kombinationen aus Testframework, Test Runner und Test Reporter erlaubt. Dass dies auch zu Unterschieden in der Testausgabe führen kann, ist in \lstref{JAVACompare} dargestellt.

\begin{lstlisting}[caption={direkter Vergleich von Junit XML, erzeugt durch Junit5 mit Jupiter in Java. Jeder Ausgabe liegen dieselben beiden Testdateien zu Grunde.\\ ''\textellipsis{}'' bezeichnet gekürzte Stellen.}, label={JAVACompare}]
(*@\underline{\texttt{Junit5, Jupiter:}}@*)
<?xml version="1.0" encoding="UTF-8"?>
<testsuite name="JUnit Jupiter" tests="5" skipped="1" failures="1" errors="0" time="0.036" hostname="DESKTOP-I9APAL0" timestamp="2026-07-04T00:45:01">
<properties>
<property name="file.encoding" value="UTF-8"/>
<property name="file.separator" value="/"/>
...
<testcase name="test_neg()" classname="BasicTest" time="0.006">
(*@\underline{\texttt{Junit5, Jupiter mit Maven surefire:}}@*)
<?xml version="1.0" encoding="UTF-8"?>
<testsuite xmlns:xsi="http://www.w3.org/2001/XMLSchema-instance" xsi:noNamespaceSchemaLocation=
"https://maven.apache.org/surefire/.../surefire-test-report.xsd" version="3.0.2" name="BasicTest" time="0.013" tests="4" errors="0" skipped="1" failures="1">
  <properties>
    <property name="java.specification.version" value="21"/>
    <property name="sun.jnu.encoding" value="UTF-8"/>
    ...
  <testcase name="test_neg" classname="BasicTest" time="0.0"/>
  ...
(*@\underline{\texttt{Junit5, Jupiter mit Gradle test task:}}@*)
  <?xml version="1.0" encoding="UTF-8"?>
<testsuite name="BasicTest" tests="4" skipped="1" failures="1" errors="0" timestamp="2026-07-03T22:47:00.701Z" hostname="DESKTOP-I9APAL0" time="0.009">
  <properties/>
  <testcase name="test_neg()" classname="BasicTest" time="0.001"/>
  ...
\end{lstlisting}
Obwohl für jede Ausgabe JUnit5 mit Jupiter als Testframework benutzt wird, jede Ausgabe JUnit XML als Format aufweist, und jede Ausgabe auf denselben beiden Testdateien basiert, gibt es durch die unterscheidlichen Test Runner und Test Reporter deutliche Unterschiede in der Testausgabe.\\
So ignoriert Gradle beispielsweise in diesem exemplarischen Fall das Element \texttt{<properties>}, während Maven und Jupiter (ohne build-tool) eine exessive Liste von \texttt{<property>} in der Testausgabe mitgeben. \\
Auch die Benennung der Testcases unterscheidet sich:\\
Jupiter und Gradle behalten den Namen \texttt{test\_neg()} bei, so wie er auch in der Testdatei verwendet wird, während Maven den Namen zu \texttt{test\_neg} kürzt.\\
Ähnliches bei der Benennung der Testsuite: Da Gradle und Maven für jede Datei eine eigene Ausgabedatei erzeugen, benutzen die beiden build-tool den Namen der jeweiligen Testsuite, hier ''BasicTest'', als Wert des Attributes ''name'' im Element \texttt{<testsuite>}. Jupiter gibt dagegen nur eine Ausgabedatei aus, die das Ergebnis aller Testdateien enthält, und benutzt in dieser den Standardwert ''Junit Jupiter'' als Namen für \texttt{<testsuite>}.

Es ist somit ersichtlich, dass die konkrete Struktur einer Testausgabe nicht allein durch das verwendete Testframework oder das gewählte Testausgabeformat bestimmt wird. Vielmehr entsteht diese durch das Zusammenwirken in diesem Kapitel beschriebenen Komponenten. Selbst bei identischem Testframework, Ausgabeformat und zugrunde liegenden Testdateien können sich die erzeugten Ausgaben hinsichtlich ihrer Struktur, ihrer Benennung und der enthaltenen Informationen unterscheiden.

Dieser Sachverhalt muss in \chapref{Vergleich Testausgabefomate} bei der Bewertung der Struktur der einzelnen Testausgebeformate berücksichtigt werden. 
